\documentclass{article}
\usepackage[utf8]{inputenc}
\usepackage[T1]{fontenc}
\usepackage{geometry}
\geometry{
    a4paper,
    left=2.5cm,
    right=2.5cm,
    top=2.5cm,
    bottom=2.5cm
}
\usepackage{graphicx}
\usepackage{float}
\usepackage{hyperref}
\usepackage[english]{babel}
\usepackage{amsmath}
\usepackage{amssymb}
\usepackage{array}
\usepackage{booktabs}
\usepackage{caption}
\usepackage{subcaption}
\usepackage{enumitem}
\usepackage{fancyhdr}
\usepackage{titlesec}
\usepackage{url}
\usepackage{textcomp}
\usepackage{xcolor}
\usepackage{listings}
\usepackage{multirow}
\usepackage{longtable}
\usepackage{tabularx}
\usepackage{times}

% Set up page style
\pagestyle{fancy}
\fancyhf{}
\rfoot{\thepage}
\renewcommand{\headrulewidth}{0pt}

% Title formatting
\titleformat{\chapter}[display]{\normalfont\bfseries}{\chaptertitle\thechapter}{0pt}{\Huge}
\titlespacing*{\chapter}{0pt}{50pt}{40pt}

% Document info
\title{Template Document}
\author{Academic Support Agent}
\date{\today}

\begin{document}
\maketitle

\vspace{-2cm}
%%%%%%%%%%%%%%%%%%%%%%%%%%%%%%%%%%%%%%%%%%%%%%%%%%%%%%%%%%%%%%%%%%%%%%%%%%%%%%%%%%%%%%%%%

\begin{titlepage}
	\centering
    %\vspace*{0.1 cm}
    \includegraphics[scale = 0.2]{logo.png}\\[1.0 cm]	% University Logo
    \textsc{\LARGE Indian Institute of Management\\ \vspace{0.5cm} Indore}\\[1.5 cm]	% University Name
	\textsc{\Large LSE 3}\\[0.5 cm]				% Course Code
	\textsc{\large Language Skills in English}\\[0.5 cm]				% Course Name
	\rule{\linewidth}{0.2 mm} \\[0.4 cm]
	{ \huge \bfseries \thetitle}\\
	\rule{\linewidth}{0.2 mm} \\[1.5 cm]
	
    {\large \thedate}\\[1 cm]
    
	\begin{minipage}{0.4\textwidth}
		\begin{flushleft} \large
			\emph{Author:}\\
			\theauthor
			\end{flushleft}
			\end{minipage}~
			\begin{minipage}{0.4\textwidth}
			\begin{flushright} \large
			\emph{Student ID:} \\
			2016IPM043									% Your Student Number
		\end{flushright}
	\end{minipage}\\
	
	
 
	\vfill
	
\end{titlepage}

%%%%%%%%%%%%%%%%%%%%%%%%%%%%%%%%%%%%%%%%%%%%%%%%%%%%%%%%%%%%%%%%%%%%%%%%%%%%%%%%%%%%%%%%%

%
%\pagebreak

\flushbottom

%%%%%%%%%%%%%%%%%%%%%%%%%%%%%%%%%%%%%%%%%%%%%%%%%%%%%%%%%%%%%%%%%%%%%%%%%%%%%%%%%%%%%%%%%

\section{Introduction}
Social media has a sudden usage upsurge these days. In 1990s, when we had the dot-com bubble inflation and every business was moving online and even the new ones coming up online, somewhere the experts knew that the brick and mortar shops which are the heart of today's trade would take the backseat when these online giants start their moves. However, they didn't anticipate that this would also mean more interaction for the customers with each other. Almost everyone who is following media these days can reiterate the facts how owner's personal social media handles - forget company's social media handles - are the new faces for interaction with the public. Many researches have been done to anticipate what makes a content popular online and shapes the public opinion. Researches have been in fields ranging from sociology to artificial intelligence. For works in sociology and artificial intelligence, see \cite{zeng2010} and \cite{wang2007}, respectively. These studies, however, are so wide in their treatment and methods that it is fairly difficult for any researcher who wants to start in this field to go through all of them. There are virtually thousands of online mediums to interact on \citep{harris2009}. This literature review tries to make it easier for a reviewer starting in this field and give them a sound idea of the previous works and their application to businesses.

\section{Works Reviewed}
We started with the works that were available freely in the public domain. Many of such social media studies are exclusive to the publishers and consultants who don't want to share the details of the case due to various anonymity clauses in their works. So, there was a limited work available to us and most of it hadn't actually been tested in market but was supported only by lab experiments/theories.

\paragraph{Literature Search Strategy}
Relevant reports and articles were searched through various mediums, including, but not limited to, \href{https://scholar.google.co.in/}{Google Scholar} for general business articles, \href{http://ieeexplore.ieee.org/Xplore/home.jsp}{IEEE Xplore} for articles based on artificial intelligence and social media, and \href{https://www.jstor.org/}{JSTOR}. We used various keywords to search: for e.g. ``business", ``social media", ``online branding", etc. These keywords were used for blanket search and we didn't made any changes to our preference for year of publication, journal, author, etc. due to lack of expert knowledge in the subject. 

\paragraph{Selection of Articles for Review}
Of the many articles that were available at our disposal, we selected a few of them for detailed review in the matter. The selections could be attributed to familiarity our topic and applicability with businesses. Social media, as a tool, is very popular with people but, in this review, we chose to focus on the methods and knowledge on how business use them to connect with people. To select the articles, we also valued the number of citation the document had. \href{https://scholar.google.co.in/}{Google Scholar} offers a quick summary on how many times the article has been cited in other articles and we factored that for selection.

\section{Literature Review}
Information flow over time has become more and more streamlined. It has been recognised long back how important communication over computers and computer-aided systems has become. It is a general topic for discussion amongst people for judging personality of person through their online activities \citep{walther2009}. It is common for people, for which many authors cite anecdotal evidences, that customers often discharge their queries and wrath on company's social media handles in public seeing which businesses have to answer them to save their public image and public perception \citep{mangold2009}.

Focusing on individuals in social media, numerous studies have been done. For instance, \cite{zhao2008} explained how people in modern day not only seek to manage impression face to face but also their personal branding online. Social media gives businesses opportunities to connect with customers and potential customers. Studies such as \cite{madera2012} and \cite{bohnert2010} also show that it is increasing it's presence in businesses manifold and is expected to expand further. Businesses often draw their impression online from how experience of users shared with the world. With the advent of technology and increasing internet penetration, more and more people can share their experiences online which if not handled well, can be a nightmare for the companies. \citep{schwartz2000}

In the job sectors of the businesses the scrutiny is further elated as it is increasingly becoming a common practice to have a thorough diligence of an employee's profile before employing her/him \citep{dekay2009}. This also means that the businesses are hiring people who can work in these environments without faking (and keeping bogus personality online image) and bringing the company to a trouble \citep{westerman2008}. When companies start moving their businesses and complaint centers online, alongwith the gain to have more free visibility, they are also subject to bigger field where the consumers are the king and the drivers of the advertisement \citep{schwartz2000}. 

This is unlike the old times when people paid regular attention to newspapers, called customer care executives for their problem, stood in queue to get them resolved \citep{culnan2010}. All of this went wrong, \cite{holt2016} argues, when the crowd on Internet became the ``cultural innovators". \cite{holt2016} further adds that this is an acute case of \textit{amplified sub-culture} which is difficult for companies to manage. The sub-cultures have varied interests and it is very difficult for companies or their consultants to come up with a \textit{one-size-fits-all} solution to the problem. ``On social media, what works for Shakira backfires for Crest and Clorox", says \cite{holt2016}. 

\cite{culnan2010} write in detail about how business can expand their reach using their Twitter and Facebook handles. They suggest campaigns that are focussed on ``call for action" from the readers which attracts and engages them. They also analysed involvement of of Fortune 500 companies in four major social medias - Twitter, Facebook, blogs and their client-hosted forums. They also suggested methods on how to choose between the various available options for social networks on the basis of user engagement and accessibility. \cite{hanna2011} further question myths and traditional knowledge on the subject and disseminates further insights on how companies can empower the users so that they can be active participants - thus increase the value of company as a whole. 

\cite{lipsman2012} outline several methods that brands try to reach their potential and present customers. They take example of Facebook on because of its popularity in USA. According to them, customers and visitors on this network can be classified in two categories - fans of brands which want to market themselves and friends of fans of those brands. According to what they say, these platforms can turn the friends of their fans into their customers. They measure the impact of each like in posts and then suggest companies on how to increase their reach going this way. \cite{ashley2015} did an empirical study with the help of two coders on 28 different businesses. They found that some of the 28 players had multiple options to keep the customers engaged. These ranged from blogs and wikis to virtual worlds and social games. They also observed that the most popular way to keep the customer engaged was through micro-blogging sites such as Twitter. The most common strategy was through ``functional appeal" where interaction was based on functionality of product or service provided. However, there is always an element of doubt so as to how well can any of the suggest methods be replicated for the reason \textit{one-size-fits-all} as \cite{holt2016} gave.

\section{Conclusion}
Social media is the new elephant which every business tries to leverage on to get themselves free reach to customers. But this reach doesn't come easy as the strategies come with various risks involved on how to keep so many people under constant eye as even slightest mistake can result in grave damages. These papers and articles fairly represent spectrum of ideas that are avaliable in academic world on how should businesses use social media sites for their advantages, when can these turn their back and become a disadvantage to the businesses themselves.

\section*{Acknowledgements}
I would like to deeply and sincerely thank \href{https://www.iimidr.ac.in/faculty/full-time-faculty/profile/?id=92}{Prof Urjani Chakravarty} for her guidance throughout the course without which it would have been very difficult to understand literature reviews in such detail. I'd also like to thanks my friends who reviewed this document and suggest improvements.

%\hspace{1 cm}--- Linus

\newpage
\end{document}